\documentclass[11pt,a4paper]{article}
\usepackage[utf8]{inputenc}
\usepackage[T1]{fontenc}
\usepackage[ngerman]{babel}
\usepackage{helvet}
\renewcommand{\familydefault}{\sfdefault}
\usepackage[left=2cm, right=2cm, top=2cm, bottom=2cm]{geometry}
\usepackage{amsmath,amssymb}
\usepackage{array}
\usepackage{booktabs}
\usepackage{tabularx}
\usepackage{tikz}
\usepackage{pgfplots}
\pgfplotsset{compat=1.18}
\usepackage{fancyhdr}
\usepackage{siunitx}
\sisetup{locale = DE}
\usepackage{xcolor}
\usepackage{marginnote}
\reversemarginpar
\usepackage{lastpage}

\pagestyle{fancy}
\fancyhf{}
\fancyhead[L]{\small Physik Klasse 10E}
\fancyhead[R]{\small Klausur Kinematik}
\fancyfoot[C]{\small Seite \thepage\ von \pageref{LastPage}}
\renewcommand{\headrulewidth}{0.4pt}

\setlength{\parindent}{0pt}
\setlength{\fboxsep}{8pt}

% Farben
\definecolor{afbgray}{RGB}{140,140,140}
\definecolor{zeitgray}{RGB}{120,120,120}

% BE + Zeit + AFB rechtsbündig
\newcommand{\BEZEIT}[3]{%
\hfill{\small\color{zeitgray}(#1\,Min, AFB\,#2)}\quad\textbf{(\_\_/#3\,BE)}%
}

\begin{document}

% === KOPFBEREICH ===
\noindent\fbox{\parbox{\dimexpr\textwidth-2\fboxsep-2\fboxrule}{
\centering
{\Large\bfseries Klausur: Kinematik}\\[0.3cm]
{\normalsize Klasse 10E \quad|\quad 90 Minuten \quad|\quad 45 BE}
}}

\vspace{0.4cm}

\begin{tabular}{@{}p{6cm}p{8cm}@{}}
Name: & \hrulefill \\[0.3cm]
Datum: & \hrulefill \\
\end{tabular}

\vspace{0.3cm}

% === NOTENSPIEGEL (15-NP Sek II MV) ===
\begin{center}
\small
\begin{tabular}{|c|c|c|c|c|c|c|c|c|c|c|c|c|c|c|c|}
\hline
\textbf{NP} & 15 & 14 & 13 & 12 & 11 & 10 & 9 & 8 & 7 & 6 & 5 & 4 & 3 & 2 & 1 \\
\hline
\textbf{ab BE} & 43 & 40,5 & 38 & 36 & 33,5 & 31,5 & 29 & 27 & 24,5 & 22,5 & 20 & 18 & 15 & 12 & 9 \\
\hline
\textbf{ab \%} & 95 & 90 & 85 & 80 & 75 & 70 & 65 & 60 & 55 & 50 & 45 & 40 & 33 & 27 & 20 \\
\hline
\end{tabular}
\end{center}

\vspace{0.3cm}

\begin{center}
\fbox{\parbox{0.95\textwidth}{
\small
\textbf{Hinweise:}
\begin{itemize}
    \item Erlaubte Hilfsmittel: Tafelwerk, Geodreieck, Lineal, Bleistift
    \item \textbf{Taschenrechner ist NICHT zugelassen}
    \item Fallbeschleunigung: $g \approx \SI{10}{\metre\per\second\squared}$
    \item Notiere bei Rechenaufgaben: Gegeben, Gesucht, Formel, Einsetzen, Ergebnis
    \item \textbf{Zeitmanagement:} Die Zeitangaben am Rand sind Empfehlungen. Wenn Du bei einer Aufgabe nicht weiterkommst, überspringe sie zunächst und bearbeite die nächste. Du kannst später zurückkehren.
\end{itemize}
}}
\end{center}

\vspace{0.5cm}

% ═══════════════════════════════════════════════════════════════
% AUFGABE 1: Grundwissen Kinematik (8 BE, AFB I)
% ═══════════════════════════════════════════════════════════════

\noindent\fbox{\parbox{\dimexpr\textwidth-2\fboxsep-2\fboxrule}{%
\textbf{Aufgabe 1: Grundwissen Kinematik}\hfill{\small\color{zeitgray}(ca.\ 9\,Min, AFB\,I)}\quad\textbf{8 BE}%
}}

\vspace{0.4cm}

\begin{enumerate}
\item[a)] \textbf{Gib an}, welche physikalische Größe die Steigung im $s$-$t$-Diagramm darstellt.

\hfill{\small\color{zeitgray}(1\,Min, AFB\,I)}\quad\textbf{(\_\_/1\,BE)}

\vspace{1.2cm}

\item[b)] \textbf{Gib an}, welche physikalische Größe die Fläche unter einem $v$-$t$-Graphen darstellt.

\hfill{\small\color{zeitgray}(1\,Min, AFB\,I)}\quad\textbf{(\_\_/1\,BE)}

\vspace{1.2cm}

\item[c)] \textbf{Rechne um:} \BEZEIT{2}{I}{2}

\vspace{0.3cm}
\hspace{1cm} 90 km/h = \underline{\hspace{3cm}} m/s

\vspace{0.5cm}
\hspace{1cm} 54 km/h = \underline{\hspace{3cm}} m/s

\end{enumerate}

\newpage

\begin{enumerate}
\item[d)] \textbf{Ordne zu:} Verbinde jede Bewegungsart (links) mit der passenden Graphenform (rechts), indem Du mit Lineal und Bleistift \textbf{saubere gerade Linien} zwischen den zusammengehörigen Kästen ziehst. \BEZEIT{5}{I}{4}
\end{enumerate}

\vspace{0.8cm}

\newcommand{\zuordnungskasten}[1]{%
\fbox{\parbox{5.5cm}{\centering\strut #1\strut}}%
}

\begin{center}
\begin{minipage}{0.45\textwidth}
\centering
\textbf{Bewegungsarten}\\[0.5cm]

\zuordnungskasten{Gleichförmige Bewegung\\im $s$-$t$-Diagramm}

\vspace{0.8cm}

\zuordnungskasten{Abbremsen\\im $v$-$t$-Diagramm}

\vspace{0.8cm}

\zuordnungskasten{Gleichförmige Bewegung\\im $v$-$t$-Diagramm}

\vspace{0.8cm}

\zuordnungskasten{Beschleunigte Bewegung\\im $s$-$t$-Diagramm}

\end{minipage}%
\hspace{1.5cm}%
\begin{minipage}{0.45\textwidth}
\centering
\textbf{Graphenformen}\\[0.5cm]

\zuordnungskasten{Waagerechte Linie}

\vspace{0.8cm}

\zuordnungskasten{Nach oben geöffnete Parabel}

\vspace{0.8cm}

\zuordnungskasten{Fallende Gerade}

\vspace{0.8cm}

\zuordnungskasten{Ansteigende Gerade}

\end{minipage}
\end{center}

\newpage

% ═══════════════════════════════════════════════════════════════
% AUFGABE 2: Freier Fall (8 BE, AFB I/II/III)
% ═══════════════════════════════════════════════════════════════

\noindent\fbox{\parbox{\dimexpr\textwidth-2\fboxsep-2\fboxrule}{%
\textbf{Aufgabe 2: Freier Fall -- Reaktionszeit}\hfill{\small\color{zeitgray}(ca.\ 9\,Min, AFB\,I/II/III)}\quad\textbf{6 BE}%
}}

\vspace{0.4cm}

Im Physikunterricht wird die Reaktionszeit mit einem fallenden Lineal gemessen. Person A lässt das Lineal ohne Vorwarnung los, Person B fängt es so schnell wie möglich.

\vspace{0.3cm}

\begin{enumerate}
\item[a)] Für den freien Fall gilt die Weg-Zeit-Beziehung $h = \frac{1}{2}g \cdot t^2$.

\textbf{Leite} aus dieser Formel eine Gleichung her, mit der Du die Fallzeit $t$ aus der Fallhöhe $h$ berechnen kannst. \BEZEIT{3}{I}{2}

\vspace{4cm}

\item[b)] \textbf{Berechne} die Reaktionszeit, wenn das Lineal $h = \SI{45}{\centi\metre}$ fällt.

\textit{Arbeitsgleichung:} $t = \sqrt{\dfrac{2h}{g}}$ \BEZEIT{3}{II}{2}

\vspace{5cm}

\item[c)] \textbf{Erkläre}, warum dieses Experiment auf dem Mond ($g_{\text{Mond}} \approx \SI{1,6}{\metre\per\second\squared}$) ungenauer wäre als auf der Erde. \BEZEIT{3}{III}{2}

\vspace{4cm}

\end{enumerate}

\newpage

% ═══════════════════════════════════════════════════════════════
% AUFGABE 3: Bremsvorgang (8 BE, AFB I/II/III)
% ═══════════════════════════════════════════════════════════════

\noindent\fbox{\parbox{\dimexpr\textwidth-2\fboxsep-2\fboxrule}{%
\textbf{Aufgabe 3: Bremsvorgang}\hfill{\small\color{zeitgray}(ca.\ 11\,Min, AFB\,I/II/III)}\quad\textbf{8 BE}%
}}

\vspace{0.4cm}

Ein PKW fährt mit \SI{108}{\kilo\metre\per\hour} auf einer Landstraße. Der Fahrer erkennt ein Hindernis und bremst sofort mit $a = \SI{-6}{\metre\per\second\squared}$.

\vspace{0.3cm}

\begin{enumerate}
\item[a)] \textbf{Rechne} die Geschwindigkeit in m/s \textbf{um}. \BEZEIT{1}{I}{1}

\vspace{2cm}

\item[b)] \textbf{Berechne} die Bremszeit bis zum Stillstand. \BEZEIT{3}{II}{2}

\vspace{4cm}

\item[c)] \textbf{Berechne} den Bremsweg. \BEZEIT{3}{II}{2}

\vspace{4cm}

\item[d)] Das Hindernis befindet sich \SI{80}{\metre} vor dem Fahrzeug. \textbf{Beurteile}, ob der Fahrer rechtzeitig stoppen kann. \textbf{Begründe} Deine Antwort mit Bezug auf Deine Berechnungen. \BEZEIT{4}{III}{3}

\vspace{4cm}

\end{enumerate}

\newpage

% ═══════════════════════════════════════════════════════════════
% AUFGABE 4: Diagrammanalyse und -konstruktion (12 BE, AFB II/III)
% ═══════════════════════════════════════════════════════════════

\noindent\fbox{\parbox{\dimexpr\textwidth-2\fboxsep-2\fboxrule}{%
\textbf{Aufgabe 4: Busfahrt durch Rostock}\hfill{\small\color{zeitgray}(ca.\ 25\,Min, AFB\,II/III)}\quad\textbf{12 BE}%
}}

\vspace{0.4cm}

Ein Linienbus fährt von der Haltestelle "`Stadtmitte"' zur Haltestelle "`Hauptbahnhof"'. Der erste Streckenabschnitt führt durch eine 50-km/h-Zone, der zweite durch eine verkehrsberuhigte 30-km/h-Zone. Das folgende $v$-$t$-Diagramm zeigt den Geschwindigkeitsverlauf.

\vspace{0.3cm}

\begin{center}
\begin{tikzpicture}
\begin{axis}[
    width=14cm, height=4.5cm,
    xlabel={$t$ in s},
    ylabel={$v$ in m/s},
    xmin=0, xmax=28,
    ymin=0, ymax=14,
    grid=both,
    grid style={gray!30},
    major grid style={gray!50},
    axis lines=left,
    xtick={0,4,10,14,18,22,28},
    ytick={0,4,8,12},
    minor tick num=1,
    every axis label/.style={font=\small},
]
% Phase 1: Anfahren (0-4s), 0→12 m/s
\addplot[very thick, black] coordinates {(0,0) (4,12)};
% Phase 2: Gleichförmig (4-10s), 12 m/s
\addplot[very thick, black] coordinates {(4,12) (10,12)};
% Phase 3: Bremsen (10-14s), 12→0 m/s
\addplot[very thick, black] coordinates {(10,12) (14,0)};
% Phase 4: Halt (14-18s), 0 m/s
\addplot[very thick, black] coordinates {(14,0) (18,0)};
% Phase 5: Anfahren (18-22s), 0→8 m/s
\addplot[very thick, black] coordinates {(18,0) (22,8)};
% Phase 6: Gleichförmig (22-28s), 8 m/s
\addplot[very thick, black] coordinates {(22,8) (28,8)};

% Phasennummern
\node at (axis cs:2,13) {\small\textbf{1}};
\node at (axis cs:7,13) {\small\textbf{2}};
\node at (axis cs:12,13) {\small\textbf{3}};
\node at (axis cs:16,1.5) {\small\textbf{4}};
\node at (axis cs:20,10) {\small\textbf{5}};
\node at (axis cs:25,10) {\small\textbf{6}};
\end{axis}
\end{tikzpicture}
\end{center}

\vspace{0.3cm}

\begin{enumerate}
\item[a)] \textbf{Beschreibe} die Bewegung des Busses in den sechs Phasen. \BEZEIT{4}{II}{3}

\vspace{0.3cm}
\begin{tabular}{@{}p{1.5cm}p{13cm}@{}}
Phase 1: & \hrulefill \\[0.4cm]
Phase 2: & \hrulefill \\[0.4cm]
Phase 3: & \hrulefill \\[0.4cm]
Phase 4: & \hrulefill \\[0.4cm]
Phase 5: & \hrulefill \\[0.4cm]
Phase 6: & \hrulefill \\
\end{tabular}

\item[b)] \textbf{Berechne} die Beschleunigung in Phase 1 (Anfahren). \BEZEIT{3}{II}{2}

\vspace{2.5cm}

\end{enumerate}

\newpage

\begin{enumerate}
\item[c)] \textbf{Ermittle} die zurückgelegte Strecke nach jeder Phase. Trage Deine Ergebnisse in die Tabelle ein. \BEZEIT{8}{II}{3}

\vspace{0.2cm}

\begin{center}
\begin{tabularx}{0.95\textwidth}{|c|c|X|c|c|}
\hline
\textbf{Phase} & \textbf{Zeit} & \textbf{Wegberechnung} {\footnotesize(z.\,B. Flächenansatz)} & $\boldsymbol{\Delta s}$ & $\boldsymbol{s_{\textbf{ges}}}$ \\
\hline
1 & 0--4\,s & & & \\[1.05cm]
\hline
2 & 4--10\,s & & & \\[1.05cm]
\hline
3 & 10--14\,s & & & \\[1.05cm]
\hline
4 & 14--18\,s & & & \\[1.05cm]
\hline
5 & 18--22\,s & & & \\[1.05cm]
\hline
6 & 22--28\,s & & & \\[1.05cm]
\hline
\end{tabularx}
\end{center}

\vspace{0.3cm}

\item[d)] \textbf{Zeichne} das zugehörige $s$-$t$-Diagramm. Beschrifte die $s$-Achse und achte auf die korrekte Kurvenform. \BEZEIT{10}{III}{4}\nopagebreak

\vspace{0.2cm}\nopagebreak

% v-t-Diagramm (Referenz)
\begin{center}
\begin{tikzpicture}
\begin{axis}[
    width=14cm, height=4cm,
    xlabel={$t$ in s},
    ylabel={$v$ in m/s},
    xmin=0, xmax=28,
    ymin=0, ymax=14,
    grid=both,
    minor tick num=1,
    grid style={gray!20},
    major grid style={gray!35},
    axis lines=left,
    xtick={0,2,4,6,8,10,12,14,16,18,20,22,24,26,28},
    ytick={0,4,8,12},
    every axis label/.style={font=\footnotesize},
    tick label style={font=\footnotesize},
]
% Kurve
\addplot[very thick, black] coordinates {(0,0) (4,12)};
\addplot[very thick, black] coordinates {(4,12) (10,12)};
\addplot[very thick, black] coordinates {(10,12) (14,0)};
\addplot[very thick, black] coordinates {(14,0) (18,0)};
\addplot[very thick, black] coordinates {(18,0) (22,8)};
\addplot[very thick, black] coordinates {(22,8) (28,8)};
% Gestrichelte vertikale Phasengrenzen
\draw[gray!70, dashed, thick] (axis cs:4,0) -- (axis cs:4,14);
\draw[gray!70, dashed, thick] (axis cs:10,0) -- (axis cs:10,14);
\draw[gray!70, dashed, thick] (axis cs:14,0) -- (axis cs:14,14);
\draw[gray!70, dashed, thick] (axis cs:18,0) -- (axis cs:18,14);
\draw[gray!70, dashed, thick] (axis cs:22,0) -- (axis cs:22,14);
% Phasennummern
\node at (axis cs:2,13) {\small\textbf{1}};
\node at (axis cs:7,13) {\small\textbf{2}};
\node at (axis cs:12,13) {\small\textbf{3}};
\node at (axis cs:16,1.5) {\small\textbf{4}};
\node at (axis cs:20,9.5) {\small\textbf{5}};
\node at (axis cs:25,9.5) {\small\textbf{6}};
\end{axis}
\end{tikzpicture}
\end{center}

\vspace{-0.3cm}

% s-t-Diagramm (leer, kariert, mit Zeitachse)
\begin{center}
\begin{tikzpicture}
\begin{axis}[
    width=14cm, height=7.2cm,
    xlabel={$t$ in s},
    ylabel={\phantom{$v$ in m/s}},
    xmin=0, xmax=28,
    ymin=0, ymax=200,
    grid=both,
    minor tick num=1,
    grid style={gray!20},
    major grid style={gray!35},
    axis lines=left,
    xtick={0,2,4,6,8,10,12,14,16,18,20,22,24,26,28},
    ytick={0,40,80,120,160,200},
    yticklabels={\phantom{12},\phantom{12},\phantom{12},\phantom{12},\phantom{12},\phantom{12}},
    every axis label/.style={font=\footnotesize},
    tick label style={font=\footnotesize},
]
% Gestrichelte vertikale Phasengrenzen
\draw[gray!70, dashed, thick] (axis cs:4,0) -- (axis cs:4,200);
\draw[gray!70, dashed, thick] (axis cs:10,0) -- (axis cs:10,200);
\draw[gray!70, dashed, thick] (axis cs:14,0) -- (axis cs:14,200);
\draw[gray!70, dashed, thick] (axis cs:18,0) -- (axis cs:18,200);
\draw[gray!70, dashed, thick] (axis cs:22,0) -- (axis cs:22,200);
\end{axis}
\end{tikzpicture}
\end{center}

\end{enumerate}

\newpage

% ═══════════════════════════════════════════════════════════════
% AUFGABE 5: Überholvorgang (8 BE, AFB II/III)
% ═══════════════════════════════════════════════════════════════

\noindent\fbox{\parbox{\dimexpr\textwidth-2\fboxsep-2\fboxrule}{%
\textbf{Aufgabe 5: Überholvorgang}\hfill{\small\color{zeitgray}(ca.\ 12\,Min, AFB\,I/II/III)}\quad\textbf{8 BE}%
}}

\vspace{0.4cm}

Auf einer Landstraße fährt ein Traktor mit konstant \SI{36}{\kilo\metre\per\hour}. Ein PKW nähert sich von hinten mit konstant \SI{72}{\kilo\metre\per\hour}. Der Abstand zwischen beiden Fahrzeugen beträgt anfangs \SI{200}{\metre}.

\vspace{0.3cm}

\begin{enumerate}
\item[a)] \textbf{Rechne} beide Geschwindigkeiten in m/s \textbf{um}. \BEZEIT{1}{I}{1}

\vspace{2.5cm}

\item[b)] \textbf{Bestimme} die Relativgeschwindigkeit des PKW bezüglich des Traktors.

\hfill{\small\color{zeitgray}(3\,Min, AFB\,II)}\quad\textbf{(\_\_/2\,BE)}

\vspace{3.2cm}

\item[c)] \textbf{Berechne}, nach welcher Zeit der PKW den Traktor erreicht. \BEZEIT{3}{II}{2}

\vspace{3.5cm}

\item[d)] Ein Verkehrsschild zeigt an, dass in \SI{500}{\metre} eine unübersichtliche Kurve beginnt. \textbf{Analysiere}, ob der PKW den Überholvorgang sicher abschließen kann, bevor er die Kurve erreicht. Für einen sicheren Überholvorgang muss der PKW den Traktor vollständig passieren (rechne mit zusätzlich \SI{50}{\metre} Sicherheitsabstand). \BEZEIT{5}{III}{3}

\vspace{5cm}

\end{enumerate}

\newpage

% ═══════════════════════════════════════════════════════════════
% AUFGABE 6: Fallschirmspringen (6 BE, AFB II/III)
% ═══════════════════════════════════════════════════════════════

\noindent\fbox{\parbox{\dimexpr\textwidth-2\fboxsep-2\fboxrule}{%
\textbf{Aufgabe 6: Fallschirmspringen}\hfill{\small\color{zeitgray}(ca.\ 6\,Min, AFB\,II/III)}\quad\textbf{3 BE}%
}}

\vspace{0.4cm}

Ein Fallschirmspringer springt aus \SI{4000}{\metre} Höhe. Zunächst fällt er im freien Fall, dann öffnet er seinen Fallschirm und schwebt langsam zu Boden.

\vspace{0.3cm}

\begin{enumerate}
\item[a)] \textbf{Erkläre}, warum der Springer im freien Fall zunächst immer schneller wird.

\hfill{\small\color{zeitgray}(2\,Min, AFB\,II)}\quad\textbf{(\_\_/1\,BE)}

\vspace{3.7cm}

\item[b)] Nach einigen Sekunden erreicht der Springer (ohne Fallschirm) eine nahezu konstante Fallgeschwindigkeit. \textbf{Erkläre} dieses Phänomen. \BEZEIT{2}{III}{1}

\vspace{5cm}

\item[c)] \textbf{Begründe}, warum die Grenzgeschwindigkeit mit geöffnetem Fallschirm deutlich geringer ist als ohne. \BEZEIT{2}{III}{1}

\vspace{4cm}

\end{enumerate}

\vspace{1cm}

\begin{center}
\textit{--- Ende der Klausur ---}
\end{center}

\end{document}
